\section{A sharp boundary for maximal-order torsion selection}
\label{sec:counterexample}

The torsion theorem allows the successful conductor to vary within a
finite range. It is tempting to replace that flexibility by an
exact-order conclusion with no valuation loss. The following explicit
counterexample prevents one natural such strengthening over arbitrarily
ramified coefficient rings. It is a correction to that proposed principle,
not a claim to refute a theorem in the cited literature.

Let \(G=C_{p^m}\), \(m\ge2\), with generator \(g\), and let \(\zeta\)
be a primitive \(p^m\)-th root of unity. In the group ring
\(\mathcal O_{\Q_p(\zeta)}[G]\), define
\[
 \nu=\prod_{a=1}^{p-1}\bigl([g]-\zeta^a[0]\bigr).
\]
We use additive notation for \(G\), with identity \(0\).
A character evaluates \([g]\) at its value on \(g\); the trivial
character, denoted \(1\), is the augmentation. Put
\(r_m=1/\varphi(p^m)\).

\begin{proposition}
\label{prop:maximal-order-counterexample}
The augmentation and maximal-order evaluations satisfy
\[
 v\bigl(\nu(1)\bigr)=(p-1)r_m=p^{1-m}=:t<1,
 \qquad
 v\bigl(\nu(\chi)\bigr)\ge2t
\]
for every character \(\chi\) of exact order \(p^m\), where the second
valuation may be infinite. In particular no such character satisfies
\(v(\nu(\chi))\le v(\nu(1))\).
\end{proposition}

\begin{proof}
Each \(1-\zeta^a\), \(1\le a\le p-1\), has valuation \(r_m\), giving
the augmentation formula. Write \(\chi(g)=\zeta^b\) with \(p\nmid b\).
Exactly one \(a\in\{1,\ldots,p-1\}\) satisfies \(a\equiv b\pmod p\).
For every other \(a\), \(v(\zeta^b-\zeta^a)=r_m\). For the matching
\(a\), the difference is zero or has valuation at least \(p r_m\),
because \(p\mid b-a\). Summing gives
\[
 v\bigl(\nu(\chi)\bigr)\ge (p-2)r_m+p r_m=2(p-1)r_m=2t.
\]
\end{proof}

For example, \(p=2,m=2\) gives
\(\nu=[g]-i[0]\). Its augmentation has valuation \(1/2\);
the two primitive evaluations have valuations \(+\infty\) and \(1\).
Meanwhile \(v_p(|G[p]|)=1\), so the weaker condition
\(v(\nu(1))<v_p(|G[p]|)\) would not justify the false maximal-order,
no-loss conclusion.

The threshold \(1/p^{m-1}\) appearing in the refined maximum-modulus
principle of Kriz--Nordentoft, Proposition 2.11, is exactly where this
example sits~\cite{kn}. This observation explains why dropping
ramification-sensitive thresholds is dangerous. It does not address all
their hypotheses or settle their separate questions about groups of
unbounded exponent.

There is an important quantifier limit to the example: \(t=p^{1-m}\)
tends to zero as \(m\) grows. It rules out no loss, and also any specified
additive loss \(\epsilon<t\) for this example, but by itself does not rule
out a theorem that first fixes \(\epsilon>0\) and then chooses \(m\)
large enough. The separate polynomial
\(\Phi_{p^H}(1+T)\), discussed in Section~\ref{sec:torsion-transfer},
vanishes at every root of prescribed exact order \(p^H\) while having
witness valuation \(1\). That is the appropriate obstruction to a
universal exact-order selector for all bounded-budget polynomial inputs.

The script \file{code/torsion_checks.py} verifies these valuations
exactly for \(p=2,3,5,7\) and \(m=2,3,4\). It also uses cyclotomic
resultants to illustrate why varying the chosen order permits
simultaneous avoidance for a finite family.

